\thispagestyle{empty}

\begin{center}
  {\large\bfseries UNIVERSITATEA „SAPIENTIA” DIN CLUJ-NAPOCA\par}
  {\large\bfseries Facultatea de Științe Tehnice și Umaniste din Târgu Mureș\par}
  {\large Specializarea: Calculatoare\par}
  
  \vspace{1.5cm}
  
  {\Large\bfseries FIȘA DE DISCIPLINĂ A PROIECTULUI DE DIPLOMĂ\par}
  {\Large\bfseries (FELADATKIÍRÁS)\par}
\end{center}

\vspace{1cm}

\noindent
\begin{minipage}[t]{0.50\textwidth}
  {\bfseries Coordonator științific:}\\
  Conf. univ. dr. Iclănzan David Andrei
\end{minipage}\hfill
\begin{minipage}[t]{0.45\textwidth}
  {\bfseries Candidat:}\\
  Nagy Csaba
\end{minipage}

\vspace{0.4cm}
\noindent
{\bfseries Anul absolvirii:} 2026

\vspace{0.8cm}

\begin{description}
  \item[\normalfont a)] {\bfseries Tema lucrării de licență (A dolgozat címe):}\\
    \textit{Recunoașterea automată a instrumentelor muzicale pe baza înregistrărilor audio / Automatikus hangszerfelismerés zenefelvételek alapján}
    
  \item[\normalfont b)] {\bfseries Problemele principale tratate (Főbb tárgyalt témakörök):}
    \begin{itemize}
      \item Studiul teoretizat al procesării semnalelor audio și al algoritmilor de învățare profundă (Deep Learning).
      \item Construirea unui set de date hibrid prin fuziunea a șase baze de date publice (IRMAS, NSynth, Medley-solos-DB, TinySOL, Good-sounds, ESC-50) și înregistrări proprii realizate cu microfon.
      \item Extragerea și analiza caracteristicilor spectrale: Log-Mel spectrograme, Short-Time Fourier Transform (STFT) și Mel-Frequency Cepstral Coefficients (MFCC).
      \item Proiectarea, optimizarea și antrenarea unei arhitecturi de rețea neuronală convoluțională bidimensională (2D CNN) în mediul Google Colab.
      \item Dezvoltarea unui modul de augmentare a datelor (zgomot de fond, reverberație) pentru îmbunătățirea robusteții modelului în condiții acustice reale.
      \item Validarea sistemului prin implementarea unui modul de recunoaștere în timp real a sunetelor captate la microfon.
    \end{itemize}
    
  \item[\normalfont c)] {\bfseries Desene obligatorii (Előírt rajzok, diagramok):}\\
    Schema bloc a fluxului de prelucrare a semnalului, reprezentarea grafică a arhitecturii rețelei neuronale 2D CNN, curbele de performanță (Loss/Accuracy) pe setul de antrenare și validare, matricea de confuzie obținută pe setul de test.
    
  \item[\normalfont d)] {\bfseries Softuri obligatorii (Előírt szoftverek):}
    \begin{itemize}
      \item Prelucrarea setului de date audio în Python 3.
      \item Proiectarea și antrenarea rețelei neuronale 2D CNN cu TensorFlow și Keras în Google Colab.
      \item Extragerea caracteristicilor spectrale ale sunetelor utilizând biblioteca Librosa.
      \item Dezvoltarea modulului de recunoaștere în timp real a sunetelor cu SoundDevice.
    \end{itemize}
    
  \item[\normalfont e)] {\bfseries Bibliografia recomandată (Ajánlott szakirodalom):}
    \begin{enumerate}[label={[\arabic*]}]
      \item Oppenheim, A. V., Schafer, R. W. - \textit{Digital Signal Processing}, Prentice Hall, 1975.
      \item Rabiner, L. R., Gold, B. - \textit{Theory and Application of Digital Signal Processing}, Prentice Hall, 1975.
      \item Goodfellow, I., Bengio, Y., Courville, A. - \textit{Deep Learning}, MIT Press, 2016.
      \item Librosa Core Developers - \textit{Librosa Package Documentation for Audio Analysis}, 2023.
    \end{enumerate}
    
  \item[\normalfont f)] {\bfseries Termene obligatorii de consultații (Konzultációs időpontok):}\\
    Săptămânal, conform planificării stabilite cu îndrumătorul de proiect.
    
  \item[\normalfont g)] {\bfseries Locul și durata practicii (Szakmai gyakorlat helye és időtartama):}\\
    Universitatea „Sapientia” din Cluj-Napoca, Facultatea de Științe Tehnice și Umaniste din Târgu Mureș, pe parcursul a 2 săptămâni.
\end{description}

\vspace{1.5cm}

\noindent
\begin{minipage}[t]{0.48\textwidth}
  \raggedright
  {\bfseries Coordonator științific:}\\
  \vspace{1.2cm}
  \makebox[5cm]{\dotfill}
\end{minipage}\hfill
\begin{minipage}[t]{0.48\textwidth}
  \raggedright
  {\bfseries Candidat:}\\
  \vspace{1.2cm}
  \makebox[5cm]{\dotfill}
\end{minipage}

\clearpage
